\chapter{Functions and Methods}
\label{ch:functions}

\begin{quote}
\textit{``Functions should do one thing. They should do it well. They should do it only.''}\\
---Robert C. Martin
\end{quote}

\vspace{1em}

Functions are the foundation of structured programming. In Novus, they're designed with systems programming in mind---explicit parameter passing, clear ownership semantics, and direct integration with the Amiga's calling conventions. Whether you're writing a simple utility function or implementing complex data structures, understanding Novus's function system is essential.

Functions are the fundamental building blocks of Novus programs. This chapter covers everything from basic function definitions to generic methods, closures, and compile-time evaluation. By the end, you'll understand how to write reusable, type-safe code that integrates seamlessly with the Novus type system.

\section{Function Basics}
\label{sec:function-basics}

\subsection{Defining Functions}
\label{subsec:defining-functions}

Functions are declared with the \texttt{fn} keyword:

\begin{lstlisting}
fn add(a: i32, b: i32) -> i32 {
    return a + b
}

fn greet(name: *u8) {
    // Prints the name to console
    println(name)
}
\end{lstlisting}

A function declaration consists of:
\begin{itemize}
    \item The \texttt{fn} keyword
    \item The function name (must be a valid identifier)
    \item Parameters in parentheses, each with name and type
    \item An optional return type after \texttt{->}
    \item The function body in braces
\end{itemize}

If no return type is specified, the function returns the unit type \texttt{()}.

\subsection{Implicit Returns}
\label{subsec:implicit-returns}

When a function body ends with an expression (not a statement), that expression becomes the function's return value:

\begin{lstlisting}
fn square(x: i32) -> i32 {
    x * x  // Implicit return - no return keyword needed
}

fn add(a: i32, b: i32) -> i32 {
    a + b
}
\end{lstlisting}

You can also use explicit \texttt{return} statements, which is required for early returns:

\begin{lstlisting}
fn safe_sqrt(x: i32) -> i32 {
    if x < 0 {
        return 0  // Early return requires 'return'
    }
    // ... compute sqrt ...
    result  // Implicit return at end
}
\end{lstlisting}

Both styles are valid. Use implicit returns for simple functions and explicit returns when you need early exit or want to be more explicit about control flow.

\subsection{Parameters}
\label{subsec:parameters}

Parameters are passed by value by default. Each parameter requires an explicit type:

\begin{lstlisting}
fn calculate(x: i32, y: i32, z: i32) -> i32 {
    return x + y * z
}
\end{lstlisting}

\textbf{Important:} Parameters are always immutable within the function body. To work with a mutable copy, assign to a local variable:

\begin{lstlisting}
fn process(x: i32) -> i32 {
    // x = x + 1  // Error: cannot mutate parameter
    var local = x  // Create mutable copy
    local = local + 1
    return local * 2
}
\end{lstlisting}

\subsection{Reference Parameters}
\label{subsec:reference-params}

Pass data by reference to avoid copying large structures:

\begin{lstlisting}
struct Player {
    x: i32,
    y: i32,
    health: i32,
    score: u32,
}

// Pass by immutable reference - cannot modify
fn get_position(player: &Player) -> (i32, i32) {
    return (player.x, player.y)
}

// Pass by mutable reference - can modify
fn damage(player: &var Player, amount: i32) {
    player.health = player.health - amount
    if player.health < 0 {
        player.health = 0
    }
}

fn main() -> i32 {
    var p = Player { x: 100, y: 200, health: 100, score: 0 }

    let (x, y) = get_position(&p)
    damage(&var p, 25)

    return 0
}
\end{lstlisting}

\begin{comingfromc}
In C, you use pointers for pass-by-reference and must manually track whether a function modifies data:

\begin{lstlisting}[language=C]
// C: Is player modified? Must check implementation!
void process(Player *player);
void print_info(const Player *player);
\end{lstlisting}

Novus makes mutability explicit at the call site:

\begin{lstlisting}
// Novus: Call site shows intent clearly
process(&var player)   // Will modify player
print_info(&player)    // Won't modify player
\end{lstlisting}

The key differences:
\begin{itemize}
    \item \texttt{\&T} (immutable reference) $\approx$ \texttt{const T*}
    \item \texttt{\&var T} (mutable reference) $\approx$ \texttt{T*}
    \item References are never null; use \texttt{*T} for nullable pointers
    \item The call site \texttt{\&var} makes modifications visible
\end{itemize}
\end{comingfromc}

The reference type syntax mirrors the declaration:
\begin{itemize}
    \item \texttt{\&T} --- immutable reference, pass with \texttt{\&value}
    \item \texttt{\&var T} --- mutable reference, pass with \texttt{\&var value}
\end{itemize}

\subsection{The \texttt{consuming} Keyword}
\label{subsec:consuming}

The \texttt{consuming} keyword indicates that a function takes ownership of a value, consuming it so it cannot be used afterward:

\begin{lstlisting}
fn take_ownership(consuming s: String) {
    // s is now owned by this function
    // ... (The original variable is no longer usable)
}

fn use_string() {
    var s = String::new()
    take_ownership(s)
    // s is no longer valid here - compilation error if used
}
\end{lstlisting}

This is particularly useful for builder patterns and type-state transformations:

\begin{lstlisting}
impl StringBuilder {
    // Consumes the builder and returns the final String
    pub fn build(consuming self) -> String {
        // Convert builder to String, consuming self
    }
}
\end{lstlisting}

\subsubsection{Consuming Parameters in AmigaOS FFI}

Many AmigaOS system calls use \texttt{consuming} to indicate ownership transfer. For example, \texttt{FreeMem} consumes its pointer parameter:

\begin{lstlisting}
from amiga::raw::exec import AllocMem, FreeMem
from amiga::raw::consts import MEMF_FAST

fn example() {
    let ptr = unsafe { AllocMem(1024, MEMF_FAST) }
    if ptr != null {
        // ... use ptr ...
        // ptr is consumed by FreeMem - cannot use it afterward
        unsafe { FreeMem(ptr, 1024) }
        // ptr is now invalid
    }
}
\end{lstlisting}

\subsection{Early Returns}
\label{subsec:early-returns}

Use explicit \texttt{return} for early exit from functions:

\begin{lstlisting}
fn safe_divide(a: i32, b: i32) -> i32 {
    if b == 0 {
        return 0  // Early return for error case
    }
    a / b  // Implicit return for normal case
}
\end{lstlisting}

\subsection{Returning Multiple Values}
\label{subsec:multiple-returns}

Use tuples to return multiple values:

\begin{lstlisting}
fn divmod(a: i32, b: i32) -> (i32, i32) {
    let quotient = a / b
    let remainder = a % b
    (quotient, remainder)  // Return tuple
}

fn main() -> i32 {
    let (q, r) = divmod(17, 5)
    // q = 3, r = 2
    return 0
}
\end{lstlisting}

\section{Visibility}
\label{sec:function-visibility}

\subsection{Public Functions}
\label{subsec:public-functions}

By default, functions are private to their module. Use \texttt{pub} to make them accessible from other modules:

\begin{lstlisting}
// Private - only accessible within this module
fn helper() -> i32 {
    return 42
}

// Public - accessible from other modules
pub fn api_function() -> i32 {
    return helper()
}
\end{lstlisting}

\subsection{Internal Visibility}
\label{subsec:internal-visibility}

The \texttt{internal} modifier makes a function visible within the same package but not to external code:

\begin{lstlisting}
// ... (illustrative - internal visibility modifier)
internal fn package_utility() -> i32 {
    return 100
}
\end{lstlisting}

\section{Traits}
\label{sec:traits}

Traits define shared behavior that types can implement. Understanding traits is essential before discussing generic functions, as generics often use trait bounds.

\subsection{Defining Traits}
\label{subsec:defining-traits}

A trait declares method signatures that implementing types must provide:

\begin{lstlisting}
pub trait Eq {
    fn eq(&self, other: &Self) -> bool
}
\end{lstlisting}

The \texttt{Self} type (capitalized) refers to the implementing type. In the example above, when implementing \texttt{Eq} for \texttt{Point}, \texttt{Self} becomes \texttt{Point}.

\subsection{Implementing Traits}
\label{subsec:implementing-traits}

Use \texttt{impl TraitName for Type} to implement a trait:

\begin{lstlisting}
struct Point {
    x: i32,
    y: i32,
}

impl Eq for Point {
    fn eq(&self, other: &Point) -> bool {
        return self.x == other.x && self.y == other.y
    }
}
\end{lstlisting}

\subsection{Common Standard Library Traits}
\label{subsec:common-traits}

Novus defines several important traits in \texttt{std::core}:

\begin{description}
    \item[\texttt{Drop}] Destructor called when value goes out of scope
    \item[\texttt{Clone}] Deep copy with potential allocation (returns \texttt{Option<Self>})
    \item[\texttt{Copy}] Marker trait for bitwise-copyable types
    \item[\texttt{Eq}] Equality comparison
    \item[\texttt{Ord}] Total ordering for comparison
    \item[\texttt{Hash}] Hashing for use in \texttt{HashMap}
    \item[\texttt{Display}] Format value to a \texttt{Formatter}; returns \texttt{bool} for success
    \item[\texttt{From<T>}] Type conversion from \texttt{T}
    \item[\texttt{Into<T>}] Type conversion to \texttt{T}
    \item[\texttt{Default}] Provide a default value for a type
\end{description}

Example implementing \texttt{Drop}:

\begin{lstlisting}
from amiga::raw::exec import FreeMem

impl Drop for Buffer {
    fn drop(&var self) {
        if self.size > 0 {
            unsafe { FreeMem(self.ptr, self.size) }
            self.ptr = null
            self.size = 0
        }
    }
}
\end{lstlisting}

\section{Methods and Impl Blocks}
\label{sec:methods}

Methods are functions associated with a type, defined within \texttt{impl} blocks.

\subsection{Defining Methods}
\label{subsec:defining-methods}

\begin{lstlisting}
struct Rectangle {
    width: u32,
    height: u32,
}

impl Rectangle {
    // Associated function (no self) - called as Rectangle::new()
    pub fn new(width: u32, height: u32) -> Rectangle {
        return Rectangle { width: width, height: height }
    }

    // Method with immutable self reference
    pub fn area(&self) -> u32 {
        return self.width * self.height
    }

    // Method with mutable self reference
    pub fn scale(&var self, factor: u32) {
        self.width = self.width * factor
        self.height = self.height * factor
    }

    // Method that consumes self
    pub fn into_tuple(self) -> (u32, u32) {
        return (self.width, self.height)
    }
}
\end{lstlisting}

\subsection{Self Parameter Types}
\label{subsec:self-types}

Methods can receive \texttt{self} in different ways:

\begin{center}
\begin{tabular}{lp{8cm}}
\textbf{Receiver} & \textbf{Meaning} \\
\hline
\texttt{\&self} & Immutable reference; can read but not modify \\
\texttt{\&var self} & Mutable reference; can read and modify \\
\texttt{self} & Takes ownership by value; consumes the instance \\
\texttt{consuming self} & Explicit ownership transfer (clearer intent) \\
\end{tabular}
\end{center}

Both \texttt{self} and \texttt{consuming self} have the same semantics---they consume the value. The \texttt{consuming} keyword makes the intent more explicit and is recommended for clarity.

\begin{lstlisting}
impl Resource {
    // Read-only access
    fn info(&self) -> i32 {
        return self.value
    }

    // Modify in place
    fn update(&var self, new_value: i32) {
        self.value = new_value
    }

    // Consume and transform - explicit consuming
    fn into_raw(consuming self) -> *u8 {
        let ptr = self.ptr
        // Don't run destructor - ownership transferred
        return ptr
    }
}
\end{lstlisting}

\subsection{Associated Functions}
\label{subsec:associated-functions}

Functions without a \texttt{self} parameter are called \emph{associated functions}. They're called using the type name:

\begin{lstlisting}
// ... (with Vec defined elsewhere)
impl<T> Vec<T> {
    // Associated function - no self
    // @zeroed creates a zero-initialized instance
    pub fn new() -> Vec<T> {
        return @zeroed(Vec<T>)
    }

    // Associated function with parameters
    pub fn with_capacity(cap: u32) -> Result<Vec<T>, ExecError> {
        // ... Allocate memory and return Vec
    }
}
\end{lstlisting}

Calling associated functions:

\begin{lstlisting}
// ... (with Vec and impl defined above)
var v = Vec::<i32>::new()
var v2 = Vec::<i32>::with_capacity(100)?
\end{lstlisting}

Associated functions are commonly used for constructors (\texttt{new}, \texttt{with\_capacity}) and factory methods.

\subsection{Calling Methods}
\label{subsec:calling-methods}

Methods are called using dot notation:

\begin{lstlisting}
var rect = Rectangle::new(10, 20)
let a = rect.area()        // Call method with &self
rect.scale(2)              // Call method with &var self
let (w, h) = rect.into_tuple()  // Consumes rect
\end{lstlisting}

The compiler automatically borrows \texttt{self} as needed. You don't need to write \texttt{(\&rect).area()}---the compiler inserts the borrow for you.

\begin{comingfromc}
In C, you pass structs to functions explicitly and use \texttt{->} for pointer access:

\begin{lstlisting}[language=C]
// C: Explicit struct passing
struct Rectangle *r = create_rect(10, 20);
int area = rect_area(r);     // Pass pointer explicitly
rect_scale(r, 2);            // No indication it modifies r
int w = r->width;            // Arrow for pointer access
\end{lstlisting}

Novus uses method syntax with automatic borrowing:

\begin{lstlisting}
// Novus: Method syntax, uniform dot access
var rect = Rectangle::new(10, 20)
let area = rect.area()     // Auto-borrow as &self
rect.scale(2)              // Auto-borrow as &var self
let w = rect.width         // Dot works for values too
\end{lstlisting}

Key differences:
\begin{itemize}
    \item No \texttt{->} operator---dot (\texttt{.}) works for both values and pointers
    \item Methods are called \texttt{value.method()}, not \texttt{method(\&value)}
    \item Auto-deref: \texttt{ptr.field} works like C's \texttt{ptr->field}
    \item Constructors are \texttt{Type::new()}, not separate \texttt{create\_type()} functions
\end{itemize}
\end{comingfromc}

\section{Generic Functions}
\label{sec:generics}

Generic functions work with multiple types using type parameters.

\subsection{Type Parameters}
\label{subsec:type-params}

\begin{lstlisting}
fn identity<T>(value: T) -> T {
    return value
}

fn swap<T>(a: &var T, b: &var T) {
    let temp = *a
    *a = *b
    *b = temp
}
\end{lstlisting}

Type parameters are declared in angle brackets after the function name.

\subsection{Trait Bounds}
\label{subsec:trait-bounds}

Use \texttt{where} clauses to constrain type parameters to types that implement specific traits:

\begin{lstlisting}
// ... (with Eq and Display imported)
pub fn expect_eq<T>(actual: T, expected: T, message: Str) where T: Eq + Display {
    if actual != expected {
        // ... Report failure - can use Eq and Display methods on T
    }
}
\end{lstlisting}

Multiple bounds are combined with \texttt{+}:

\begin{lstlisting}
fn sort<T>(items: &var [T], count: u32) where T: Ord + Clone {
    // Can use Ord for comparison, Clone for copying
}
\end{lstlisting}

\subsection{Turbofish Syntax}
\label{subsec:turbofish}

When the compiler cannot infer generic types, use the turbofish syntax (\texttt{::<T>}) to specify them explicitly:

\begin{lstlisting}
// Type inferred from variable declaration
let v: Vec<i32> = Vec::new()

// Explicit type with turbofish
let v = Vec::<i32>::new()

// Turbofish on function calls
let chan = bounded_channel::<i32>()?
\end{lstlisting}

The turbofish is especially useful when calling generic functions where the type cannot be inferred from context.

\subsection{Generic Trait Implementations}
\label{subsec:generic-impl}

Implement traits for generic types:

\begin{lstlisting}
// ... (with HashMap defined elsewhere)
impl<K, V> Drop for HashMap<K, V> where K: Hash + Eq {
    fn drop(&var self) {
        // ... Free all buckets and entries
    }
}
\end{lstlisting}

\section{Function Overloading}
\label{sec:overloading}

Novus supports function overloading---multiple functions with the same name but different parameter signatures.

\subsection{Overloading Rules}
\label{subsec:overload-rules}

Functions can be overloaded when they differ by:
\begin{itemize}
    \item Parameter types (\texttt{i32} vs \texttt{i64})
    \item Parameter count
    \item Reference vs value (\texttt{T} vs \texttt{\&T})
    \item Mutability (\texttt{\&T} vs \texttt{\&var T})
    \item Pointer vs value (\texttt{T} vs \texttt{*T})
\end{itemize}

\begin{lstlisting}
fn abs(x: i32) -> i32 {
    if x < 0 { return -x }
    return x
}

fn abs(x: i16) -> i16 {
    if x < 0 { return -x }
    return x
}

fn print_value(x: i32) {
    // Print i32
}

fn print_value(x: i32, y: i32) {
    // Print two values
}

fn print_value(x: i32, y: i32, z: i32) {
    // Print three values
}
\end{lstlisting}

\subsection{Overloading Restrictions}
\label{subsec:overload-restrictions}

Overloading is \textbf{not} allowed when functions differ only by:
\begin{itemize}
    \item Return type (with identical parameters)
    \item Parameter names (with identical types)
\end{itemize}

\begin{lstlisting}
// ERROR: Same parameter types, different return types
fn get() -> i32 { return 42 }
fn get() -> i64 { return 42 }  // Error: duplicate signature

// ERROR: Same types, different names
fn test(x: i32) -> i32 { return x }
fn test(y: i32) -> i32 { return y }  // Error: duplicate signature
\end{lstlisting}

\section{Const Functions}
\label{sec:const-fn}

Const functions can be evaluated at compile time when called with constant arguments.

\subsection{Defining Const Functions}
\label{subsec:const-fn-def}

\begin{lstlisting}
const fn times_two(x: i32) -> i32 {
    return x * 2
}

const fn max(a: i32, b: i32) -> i32 {
    if a > b {
        return a
    }
    return b
}

// Const fn calling another const fn
const fn clamp(value: i32, low: i32, high: i32) -> i32 {
    if value < low { return low }
    if value > high { return high }
    return value
}
\end{lstlisting}

\subsection{Using Const Functions}
\label{subsec:const-fn-use}

Const functions enable compile-time computation for constants:

\begin{lstlisting}
const WIDTH: i32 = 320
const HEIGHT: i32 = 200
const DOUBLED_WIDTH: i32 = times_two(WIDTH)    // 640 at compile time
const SCREEN_SIZE: i32 = WIDTH * HEIGHT        // 64000 at compile time
const CLAMPED: i32 = clamp(999, 0, 100)        // 100 at compile time
\end{lstlisting}

Const functions can also be called at runtime with non-constant arguments:

\begin{lstlisting}
fn process(x: i32) -> i32 {
    return times_two(x)  // Runtime call
}
\end{lstlisting}

\subsection{Const Function Restrictions}
\label{subsec:const-fn-restrictions}

Const functions have limitations:
\begin{itemize}
    \item No heap allocation
    \item No I/O operations
    \item No external function calls
    \item Only call other const functions
    \item Limited recursion depth
\end{itemize}

\section{Closures}
\label{sec:closures}

Closures are anonymous functions that can capture variables from their environment.

\subsection{Closure Syntax}
\label{subsec:closure-syntax}

\begin{lstlisting}
// Basic closure with type annotations
let add = |x: i32, y: i32| -> i32 { x + y }

// Closure capturing environment
let multiplier = 10
let scale = |x: i32| -> i32 { x * multiplier }
let result = scale(5)  // 50

// No parameters (note: space between pipes required)
let get_value = | | -> i32 { 42 }
\end{lstlisting}

Closures use pipes (\texttt{|}) to delimit parameters:

\begin{lstlisting}[language=bash]
|parameters| -> ReturnType { body }
\end{lstlisting}

\subsection{Capturing Variables}
\label{subsec:capturing}

Closures capture variables from their enclosing scope. By default, captures are by value (copied):

\begin{lstlisting}
fn example() -> i32 {
    let offset = 10
    let add_offset = |x: i32| -> i32 {
        x + offset  // Captures 'offset' by value
    }
    add_offset(5)  // 15
}
\end{lstlisting}

\section{External Functions}
\label{sec:extern-fn}

External functions declare interfaces to C code or AmigaOS system libraries.

\subsection{Extern Declarations}
\label{subsec:extern-decl}

External functions declare interfaces to AmigaOS system libraries. Here are common examples from exec.library:

\begin{lstlisting}
// Memory allocation
extern pub fn AllocMem(byteSize: u32, requirements: u32) -> *u8
extern pub fn FreeMem(consuming memoryBlock: *u8, byteSize: u32)

// Task synchronization
extern pub fn Wait(signalSet: u32) -> u32
extern pub fn Signal(task: *Task, signals: u32)

// Message passing
extern pub fn PutMsg(port: *MsgPort, message: *Message)
extern pub fn GetMsg(port: *MsgPort) -> *Message
\end{lstlisting}

From dos.library (note BCPL pointer conventions---file handles are \texttt{i32}):

\begin{lstlisting}
// File operations (using BCPL pointers)
extern pub fn Open(name: *u8, accessMode: i32) -> i32
extern pub fn Close(consuming file: i32) -> i32
extern pub fn Read(file: i32, buffer: *u8, length: i32) -> i32
extern pub fn Write(file: i32, buffer: *u8, length: i32) -> i32
\end{lstlisting}

External functions have no body---they're implemented in ROM or disk-based libraries and linked at runtime.

\subsection{Calling External Functions}
\label{subsec:calling-extern}

AmigaOS system calls require \texttt{unsafe} blocks because they bypass Novus's safety guarantees:

\begin{lstlisting}
from amiga::raw::exec import AllocMem, FreeMem
from amiga::raw::consts import MEMF_FAST, MEMF_CLEAR

fn allocate_buffer(size: u32) -> *u8 {
    var ptr: *u8 = null
    unsafe {
        ptr = AllocMem(size, MEMF_FAST | MEMF_CLEAR)
    }
    return ptr
}

fn free_buffer(ptr: *u8, size: u32) {
    unsafe {
        FreeMem(ptr, size)
    }
}
\end{lstlisting}

\subsection{AmigaOS Library Bindings}
\label{subsec:ffi-bindings}

Novus provides pre-generated FFI bindings for all AmigaOS libraries. These are auto-generated from SFD (System Function Definition) files:

\begin{lstlisting}
// Import from pre-generated FFI modules
from amiga::raw::exec import AllocMem, FreeMem, Wait, Signal
from amiga::raw::dos import Open, Close, Read, Write
from amiga::raw::graphics import BltBitMap, BltTemplate
from amiga::raw::intuition import OpenWindow, CloseWindow

// Import structure definitions
from amiga::raw::structs import Task, MsgPort, BitMap, RastPort

// Import constants
from amiga::raw::consts import MEMF_CHIP, MEMF_FAST, MODE_OLDFILE
\end{lstlisting}

Developers should import from these modules rather than writing \texttt{extern fn} declarations manually.

\subsection{Variadic Functions}
\label{subsec:variadic}

External functions can accept variable arguments using \texttt{...args}:

\begin{lstlisting}
extern fn VPrintf(format: *u8, ...args) -> i32
\end{lstlisting}

Note that variadic functions are only supported for external declarations; Novus code cannot define variadic functions.

\section{Best Practices}
\label{sec:fn-best-practices}

\subsection{Parameter Passing Guidelines}
\label{subsec:param-guidelines}

\begin{center}
\begin{tabular}{lp{8cm}}
\textbf{Situation} & \textbf{Recommendation} \\
\hline
Small values (integers, bools) & Pass by value \\
Large structs, read-only & Pass by reference (\texttt{\&T}) \\
Large structs, modify & Pass by mutable reference (\texttt{\&var T}) \\
Transfer ownership & Pass by value or use \texttt{consuming} \\
\end{tabular}
\end{center}

\subsection{Method Receiver Guidelines}
\label{subsec:receiver-guidelines}

\begin{center}
\begin{tabular}{lp{8cm}}
\textbf{Receiver} & \textbf{Use When} \\
\hline
\texttt{\&self} & Reading data, computing derived values \\
\texttt{\&var self} & Modifying internal state \\
\texttt{self} / \texttt{consuming self} & Transforming into another type, builder finalization \\
\end{tabular}
\end{center}

\subsection{AmigaOS FFI Best Practices}
\label{subsec:amiga-best-practices}

When calling AmigaOS system libraries:

\begin{itemize}
    \item Always wrap library calls in \texttt{unsafe} blocks
    \item Check return values---many functions return \texttt{NULL}/\texttt{0} on failure
    \item Track allocation sizes---\texttt{FreeMem} requires the original size
    \item Match memory flags---free with the same size you allocated
    \item Respect \texttt{consuming} parameters---don't use pointers/handles after consumption
    \item Use stdlib wrappers (\texttt{MemoryBlock}, \texttt{Vec}, etc.) when possible for safety
    \item BCPL pointers (DOS.library) are \texttt{i32} values, not \texttt{*u8}
\end{itemize}

Example of safe memory management:

\begin{lstlisting}
from std::memory::block import MemoryBlock
from amiga::raw::consts import MEMF_FAST

// Prefer MemoryBlock over raw AllocMem/FreeMem
match MemoryBlock::alloc(1024, MEMF_FAST) {
    Option::Some(var block) => {
        // Use block.ptr() to access raw pointer
        let ptr = block.ptr()
        // ... use ptr ...
        // Automatic cleanup when block goes out of scope
    },
    Option::None => {
        // Handle allocation failure
    }
}
\end{lstlisting}

\subsection{Naming Conventions}
\label{subsec:naming-conventions}

\begin{itemize}
    \item Use \texttt{snake\_case} for function names: \texttt{get\_value}, \texttt{set\_position}
    \item Use \texttt{new} for constructors: \texttt{Vec::new()}, \texttt{String::new()}
    \item Use \texttt{with\_*} for constructors with parameters: \texttt{with\_capacity}
    \item Use \texttt{is\_*} for boolean queries: \texttt{is\_empty}, \texttt{is\_valid}
    \item Use \texttt{into\_*} for consuming conversions: \texttt{into\_raw}, \texttt{into\_string}
    \item Use \texttt{as\_*} for borrowing conversions: \texttt{as\_str}, \texttt{as\_slice}
    \item Use \texttt{try\_*} for fallible operations returning \texttt{Result}: \texttt{try\_alloc}
\end{itemize}

\section{Common Function Patterns}
\label{sec:fn-patterns}

This section presents idiomatic patterns for function design in Novus.

\subsection{The Builder Pattern}

Builders accumulate configuration before producing a final value:

\begin{lstlisting}
struct WindowBuilder {
    title: *u8,
    width: u16,
    height: u16,
    left: u16,
    top: u16,
    flags: u32,
}

impl WindowBuilder {
    pub fn new(title: *u8) -> WindowBuilder {
        return WindowBuilder {
            title: title,
            width: 640,
            height: 400,
            left: 0,
            top: 0,
            flags: 0,
        }
    }

    // Each setter returns self for chaining
    pub fn size(&var self, width: u16, height: u16) -> &var WindowBuilder {
        self.width = width
        self.height = height
        return self
    }

    pub fn position(&var self, left: u16, top: u16) -> &var WindowBuilder {
        self.left = left
        self.top = top
        return self
    }

    pub fn borderless(&var self) -> &var WindowBuilder {
        self.flags |= WFLG_BORDERLESS
        return self
    }

    pub fn backdrop(&var self) -> &var WindowBuilder {
        self.flags |= WFLG_BACKDROP
        return self
    }

    // Final build consumes the builder
    pub fn build(consuming self) -> Result<*Window, IntuitionError> {
        // Create the window from accumulated settings
        let window = unsafe {
            OpenWindowTags(
                null,
                WA_Title, self.title,
                WA_Width, (u32)self.width,
                WA_Height, (u32)self.height,
                WA_Left, (u32)self.left,
                WA_Top, (u32)self.top,
                WA_Flags, self.flags,
                TAG_DONE
            )
        }

        if window == null {
            return Result::Err(IntuitionError::OpenFailed)
        }
        return Result::Ok(window)
    }
}

// Usage
fn open_game_window() -> Result<*Window, IntuitionError> {
    var builder = WindowBuilder::new("My Game")
    let window = builder
        .size(320, 256)
        .position(0, 0)
        .borderless()
        .backdrop()
        .build()?

    return Result::Ok(window)
}
\end{lstlisting}

\subsection{The Factory Pattern}

Factories create objects based on runtime conditions:

\begin{lstlisting}
enum AudioDriver {
    Paula { channels: u8 },
    Ahi { device_id: i32 },
    Null,
}

impl AudioDriver {
    // Factory method chooses implementation at runtime
    pub fn create(prefer_ahi: bool) -> Result<AudioDriver, AudioError> {
        if prefer_ahi {
            // Try AHI first
            match try_open_ahi() {
                Result::Ok(id) => {
                    return Result::Ok(AudioDriver::Ahi { device_id: id })
                },
                Result::Err(_) => {
                    // Fall through to Paula
                }
            }
        }

        // Try Paula hardware
        if is_paula_available() {
            return Result::Ok(AudioDriver::Paula { channels: 4 })
        }

        // No audio available - return null driver
        return Result::Ok(AudioDriver::Null)
    }

    pub fn play(&self, sample: &Sample) -> Result<(), AudioError> {
        match *self {
            AudioDriver::Paula { channels } => play_paula(sample, channels),
            AudioDriver::Ahi { device_id } => play_ahi(sample, device_id),
            AudioDriver::Null => Result::Ok(()),  // Silent success
        }
    }
}
\end{lstlisting}

\subsection{The Callback Pattern}

Pass functions as parameters for customizable behavior:

\begin{lstlisting}
// Type alias for callback signature
type Comparator<T> = fn(&T, &T) -> i32

// Sort with custom comparator
pub fn sort_by<T>(items: &var [T], count: u32, compare: Comparator<T>) {
    // Simple bubble sort for demonstration
    for i in 0..count {
        for j in 0..(count - 1 - i) {
            if compare(&items[j], &items[j + 1]) > 0 {
                // Swap
                let temp = items[j]
                items[j] = items[j + 1]
                items[j + 1] = temp
            }
        }
    }
}

// Usage
fn example() {
    var scores = [100, 50, 75, 25, 90]

    // Sort ascending
    sort_by(&var scores, 5, |a: &i32, b: &i32| -> i32 { *a - *b })

    // Sort descending
    sort_by(&var scores, 5, |a: &i32, b: &i32| -> i32 { *b - *a })
}
\end{lstlisting}

\subsection{The RAII Guard Pattern}

Functions that return guards for automatic cleanup:

\begin{lstlisting}
struct ForbidGuard {
    _dummy: u8,
}

impl ForbidGuard {
    pub fn new() -> ForbidGuard {
        unsafe { Forbid() }
        return ForbidGuard { _dummy: 0 }
    }
}

impl Drop for ForbidGuard {
    fn drop(&var self) {
        unsafe { Permit() }
    }
}

// Usage - interrupts disabled for the scope
fn critical_section() {
    let _guard = ForbidGuard::new()
    // Forbid() is active here
    modify_shared_data()
    // Permit() called automatically when _guard goes out of scope
}
\end{lstlisting}

\subsection{The Visitor Pattern}

Process structures without knowing their internals:

\begin{lstlisting}
trait FileVisitor {
    fn visit_file(&var self, path: &Str, size: u32)
    fn visit_directory(&var self, path: &Str)
    fn should_descend(&self, path: &Str) -> bool
}

fn walk_directory(path: &Str, visitor: &var dyn FileVisitor) -> Result<(), DosError> {
    let lock = Lock(path.as_ptr(), ACCESS_READ)
    if lock == 0 {
        return Result::Err(DosError::NotFound)
    }
    defer { UnLock(lock) }

    var fib = @zeroed(FileInfoBlock)
    if !Examine(lock, &var fib) {
        return Result::Err(DosError::ExamineFailed)
    }

    while ExNext(lock, &var fib) {
        let name = fib_name(&fib)

        if fib.fib_DirEntryType > 0 {
            // Directory
            visitor.visit_directory(&name)
            if visitor.should_descend(&name) {
                let sub_path = path.concat(&name)
                walk_directory(&sub_path, visitor)?
            }
        } else {
            // File
            visitor.visit_file(&name, fib.fib_Size as u32)
        }
    }

    return Result::Ok(())
}

// Example visitor
struct FileSizeCounter {
    total_bytes: u64,
    file_count: u32,
}

impl FileVisitor for FileSizeCounter {
    fn visit_file(&var self, path: &Str, size: u32) {
        self.total_bytes += (u64)size
        self.file_count += 1
    }

    fn visit_directory(&var self, path: &Str) {
        // Ignore directories
    }

    fn should_descend(&self, path: &Str) -> bool {
        return true  // Descend into all directories
    }
}
\end{lstlisting}

\section{Showcase: Animation System}

This showcase demonstrates functions and methods in a practical animation system:

\begin{lstlisting}
// animation_system_showcase.novus
// Demonstrates: functions, methods, traits, generics, closures

from std::io import println
from std::core import Option, Clone
from std::collections::vec import Vec

// ============================================================================
// Easing Functions
// ============================================================================

type EasingFn = fn(fixed32) -> fixed32

// Linear interpolation (no easing)
pub fn ease_linear(t: fixed32) -> fixed32 {
    return t
}

// Quadratic ease-in
pub fn ease_in_quad(t: fixed32) -> fixed32 {
    return t * t
}

// Quadratic ease-out
pub fn ease_out_quad(t: fixed32) -> fixed32 {
    return t * (2.0 - t)
}

// Quadratic ease-in-out
pub fn ease_in_out_quad(t: fixed32) -> fixed32 {
    if t < 0.5 {
        return 2.0 * t * t
    }
    return -1.0 + (4.0 - 2.0 * t) * t
}

// ============================================================================
// Animatable Trait
// ============================================================================

trait Animatable {
    fn lerp(&self, target: &Self, t: fixed32) -> Self
}

impl Animatable for i32 {
    fn lerp(&self, target: &i32, t: fixed32) -> i32 {
        let diff = *target - *self
        return *self + (fixed32::from_int(diff) * t).to_int()
    }
}

impl Animatable for fixed32 {
    fn lerp(&self, target: &fixed32, t: fixed32) -> fixed32 {
        return *self + (*target - *self) * t
    }
}

// ============================================================================
// Animation State
// ============================================================================

enum AnimationState {
    Playing,
    Paused,
    Finished,
}

// ============================================================================
// Tween Structure
// ============================================================================

struct Tween<T> where T: Animatable + Clone {
    start_value: T,
    end_value: T,
    duration_frames: u32,
    current_frame: u32,
    easing: EasingFn,
    state: AnimationState,
}

impl<T> Tween<T> where T: Animatable + Clone {
    // Constructor with default linear easing
    pub fn new(start: T, end: T, duration: u32) -> Tween<T> {
        return Tween {
            start_value: start,
            end_value: end,
            duration_frames: duration,
            current_frame: 0,
            easing: ease_linear,
            state: AnimationState::Playing,
        }
    }

    // Builder method to set easing function
    pub fn with_easing(&var self, easing: EasingFn) -> &var Tween<T> {
        self.easing = easing
        return self
    }

    // Get current interpolated value
    pub fn value(&self) -> T {
        if self.duration_frames == 0 {
            return self.end_value.clone().unwrap()
        }

        // Calculate normalized time (0.0 to 1.0)
        let t_raw = fixed32::from_int((i32)self.current_frame)
            / fixed32::from_int((i32)self.duration_frames)

        // Clamp to [0, 1]
        let t_clamped = if t_raw > 1.0 { 1.0 } else { t_raw }

        // Apply easing function
        let t_eased = (self.easing)(t_clamped)

        // Interpolate between start and end
        return self.start_value.lerp(&self.end_value, t_eased)
    }

    // Advance animation by one frame
    pub fn tick(&var self) {
        match self.state {
            AnimationState::Playing => {
                self.current_frame += 1
                if self.current_frame >= self.duration_frames {
                    self.state = AnimationState::Finished
                }
            },
            _ => { }
        }
    }

    // Check if animation is complete
    pub fn is_finished(&self) -> bool {
        match self.state {
            AnimationState::Finished => true,
            _ => false,
        }
    }

    // Pause the animation
    pub fn pause(&var self) {
        match self.state {
            AnimationState::Playing => {
                self.state = AnimationState::Paused
            },
            _ => { }
        }
    }

    // Resume the animation
    pub fn resume(&var self) {
        match self.state {
            AnimationState::Paused => {
                self.state = AnimationState::Playing
            },
            _ => { }
        }
    }

    // Reset to beginning
    pub fn reset(&var self) {
        self.current_frame = 0
        self.state = AnimationState::Playing
    }
}

// ============================================================================
// Animation Sequence
// ============================================================================

struct AnimationSequence<T> where T: Animatable + Clone {
    tweens: Vec<Tween<T>>,
    current_index: u32,
    loop_count: i32,  // -1 for infinite
    loops_remaining: i32,
}

impl<T> AnimationSequence<T> where T: Animatable + Clone {
    pub fn new() -> Option<AnimationSequence<T>> {
        let tweens = Vec::<Tween<T>>::new()
        return Option::Some(AnimationSequence {
            tweens: tweens,
            current_index: 0,
            loop_count: 1,
            loops_remaining: 1,
        })
    }

    pub fn add(&var self, tween: Tween<T>) {
        self.tweens.push(tween)
    }

    pub fn set_loop_count(&var self, count: i32) {
        self.loop_count = count
        self.loops_remaining = count
    }

    pub fn value(&self) -> Option<T> {
        if self.current_index >= self.tweens.len() {
            return Option::None
        }
        return Option::Some(self.tweens.get(self.current_index).unwrap().value())
    }

    pub fn tick(&var self) -> bool {
        if self.current_index >= self.tweens.len() {
            return false  // Sequence complete
        }

        let current = self.tweens.get_mut(self.current_index).unwrap()
        current.tick()

        if current.is_finished() {
            self.current_index += 1

            // Check if we should loop
            if self.current_index >= self.tweens.len() {
                if self.loop_count < 0 || self.loops_remaining > 1 {
                    // Reset for next loop
                    self.current_index = 0
                    for i in 0..self.tweens.len() {
                        self.tweens.get_mut(i).unwrap().reset()
                    }
                    if self.loops_remaining > 0 {
                        self.loops_remaining -= 1
                    }
                    return true  // Still playing
                }
                return false  // Done
            }
        }

        return true  // Still playing
    }
}

// ============================================================================
// Example Usage
// ============================================================================

fn main() -> i32 {
    println("=== Animation System Showcase ===\n")

    // Simple tween with linear easing
    println("--- Linear Tween (0 to 100, 10 frames) ---")
    var linear_tween = Tween::new(0, 100, 10)

    while !linear_tween.is_finished() {
        println(f"Frame {linear_tween.current_frame}: value = {linear_tween.value()}")
        linear_tween.tick()
    }

    // Tween with custom easing
    println("\n--- Ease-Out Quad Tween ---")
    var eased_tween = Tween::new(0, 100, 10)
    eased_tween.with_easing(ease_out_quad)

    while !eased_tween.is_finished() {
        println(f"Frame {eased_tween.current_frame}: value = {eased_tween.value()}")
        eased_tween.tick()
    }

    // Animation sequence
    println("\n--- Animation Sequence (3 tweens, 2 loops) ---")
    match AnimationSequence::<i32>::new() {
        Option::Some(var seq) => {
            seq.add(Tween::new(0, 50, 5))
            seq.add(Tween::new(50, 25, 5))
            seq.add(Tween::new(25, 100, 5))
            seq.set_loop_count(2)

            var frame = 0
            while seq.tick() {
                match seq.value() {
                    Option::Some(val) => {
                        println(f"Frame {frame}: value = {val}")
                    },
                    Option::None => { }
                }
                frame += 1
            }
        },
        Option::None => {
            println("Failed to create sequence")
        }
    }

    println("\n=== Showcase Complete ===")
    return 0
}
\end{lstlisting}

This showcase demonstrates:
\begin{itemize}
    \item \textbf{Function types} as easing function callbacks
    \item \textbf{Traits} for defining animatable behavior
    \item \textbf{Generic structs and methods} with trait bounds
    \item \textbf{Builder pattern} with \texttt{with\_easing}
    \item \textbf{Consuming and reference self} in methods
    \item \textbf{Match expressions} for state handling
    \item \textbf{Closures} (implicitly in function type usage)
\end{itemize}

\section{Summary}
\label{sec:fn-summary}

This chapter covered the complete function system in Novus:

\begin{itemize}
    \item \textbf{Function basics}: Declaration, parameters, return values
    \item \textbf{Reference parameters}: \texttt{\&T} for reading, \texttt{\&var T} for modifying
    \item \textbf{Traits}: Defining and implementing shared behavior
    \item \textbf{Methods}: Impl blocks, self receivers, associated functions
    \item \textbf{Generics}: Type parameters, trait bounds, turbofish syntax
    \item \textbf{Overloading}: Multiple functions with the same name
    \item \textbf{Const functions}: Compile-time evaluation
    \item \textbf{Closures}: Anonymous functions capturing environment
    \item \textbf{External functions}: AmigaOS FFI and system calls
\end{itemize}

Functions in Novus are designed for clarity and safety. Explicit parameter passing, trait bounds, and ownership semantics help prevent bugs while keeping the code readable.

The next chapter explores modules and project organization, showing how to structure larger Novus programs.
